%! Boxplots, Outliers, and Resistance — Unit 2, Lesson 2.3 — Homework (Answer Key)
% PILOT SHAPE, Unit 2 timing (5 / 45 / 5 / 5): the lesson's SCORED individual practice and the
% LAST component in the packet. ASSIGNED at Close & Assign and done OUTSIDE class (it is not
% started in class from Unit 2 on). Packet pages are the default; the plan records the call.
% THIRD CONTEXT: the notes teach on Pine Street Pizza, the Guided Practice runs on the Maple
% Grove shelter, and these items are on Elm Park Bike Shop — recall is not the skill.
% TWO PAGES MAX: two boxes — items 1-2 on page 1, items 3-6 and the spiralbox on page 2 — so
% no item breaks across a page.
% ROLES: 1 core procedure (five-number summary) · 2 fences + the "$40 looks far" check (the
% secondary misconception, justify with a fence) · 3 the CONTRAST PAIR (mean vs median, with
% and without the $900) · 4 the CRUX head-on (the classmate's claim; which summaries resist) ·
% 5 interpret-in-context (cumulative relative frequency graph) · 6 SPIRAL to Lesson 2.2 (the
% range is not resistant) + justify the spread to report. 2 points each, 12 total.
% THE DATA — Elm Park, eleven used bikes ($): 40 110 120 130 140 150 160 180 180 200 900
%   median 150 (6th); Q1 = med(40,110,120,130,140) = 120; Q3 = med(160,180,180,200,900) = 180;
%   IQR 60; step 90; fences 30 and 270; 900 the only outlier; 40 > 30 is NOT one; whisker
%   stops at 200. Sum 2310, mean 210.
%   Without the 900 (ten): sum 1410, mean 141; median (140+150)/2 = 145; Q1 = med(40..140) = 120;
%   Q3 = med(150,160,180,180,200) = 180; IQR 60. Range 860 -> 160.
%   Sample SD: sqrt(542400/10) = 232.9 -> $233 with; sqrt(18690/9) = 45.6 -> $46 without.
%   Season (40 bikes) cumulative %: $50 0, $100 15, $150 50, $200 75, $250 90, $300 95, $350 100.
% PAGE LOCKSTEP: the key differs only in saar-key, the header, \blank -> \ans, \writeline ->
% \ansline; every work block is byte-identical.
\documentclass[12pt]{article}
\usepackage{saar-article}
\usepackage{saar-key}   % pulls in -boxes; do NOT also load -boxes
\newcommand{\TallMath}[1]{$\displaystyle #1\rule[-1.4em]{0pt}{3.2em}$}
% Word bank strip for every fill-in cluster (user request 2026-09-06). Defined per document;
% the key inherits it and prints the same strip, so heights match.
\newcommand{\wordbank}[1]{\par\vspace{2pt}\noindent\fcolorbox{goldacc}{goldbg}{\parbox{\dimexpr\linewidth-2\fboxsep-2\fboxrule\relax}{\small\textbf{Word bank:} #1}}\par\vspace{4pt}}

\begin{document}

\pageheader{Unit 2, Lesson 2.3}{Homework --- Answer Key}
% NAMESTRIP: no \namedateperiod here — mirrors the blank. NO teachernote here —
% teacher prose lives in the lesson plan.

\vspace{0.04in}

% Authored identically in homework/ and homework_key/.
\begin{remindbox}
\small
\textbf{This is your graded homework.} It is scored out of 12 (2 points per item) and due at
the start of the first class after two study halls. It is assigned today and done outside
class --- the \emph{Keep in Mind} box on the cover is the page to flip back to. Use the notes'
quartile method, and say every answer about the bikes.
\end{remindbox}

\vspace{0.03in}

\boxguard[20]
\begin{scenariobox}[Elm Park Bike Shop]{cerulean}
Elm Park Bike Shop has eleven used bikes for sale. Their prices, in order, in dollars:
\[ 40 \quad 110 \quad 120 \quad 130 \quad 140 \quad 150 \quad 160 \quad 180 \quad 180 \quad 200
   \quad 900 \]
The \$900 bike is a carbon-fiber racer someone traded in; the \$40 bike needs new tires.
\end{scenariobox}

\vspace{0.03in}

\boxguard
\begin{notesbox}{Boxplots, Outliers, and Resistance}
Every answer names the bikes and the dollars.
\wordbank{outlier \;\textbullet\; fence \;\textbullet\; resistant \;\textbullet\; not resistant \;\textbullet\; mean \;\textbullet\; median \;\textbullet\; IQR \;\textbullet\; standard deviation \;\textbullet\; range}
\begin{enumerate}[leftmargin=1.6em, itemsep=5pt, topsep=2pt]

  \item Find the five-number summary of the eleven prices. Show the two halves you used for
        $Q_1$ and $Q_3$.

        \begin{center}
        \renewcommand{\arraystretch}{1.6}
        \begin{tabularx}{\linewidth}{@{} Y Y Y Y Y @{}}
        \toprule
        \textbf{Min} & $\boldsymbol{Q_1}$ & \textbf{Median} & $\boldsymbol{Q_3}$ & \textbf{Max} \\
        \midrule
        \ans{\$40} & \ans{\$120} & \ans{\$150} & \ans{\$180} & \ans{\$900} \\
        \bottomrule
        \end{tabularx}
        \end{center}

        \begin{work}
          Q_1 &= 120 \quad (\text{median of } 40, 110, 120, 130, 140) \\
          Q_3 &= 180 \quad (\text{median of } 160, 180, 180, 200, 900)
        \end{work}

  \item Use the $1.5 \times \text{IQR}$ rule to find the fences.

        \begin{work}
          \text{IQR} &= 180 - 120 = 60, \quad 1.5 \times \text{IQR} = 1.5(60) = 90 \\
          \text{lower fence} &= 120 - 90 = 30 \\
          \text{upper fence} &= 180 + 90 = 270
        \end{work}

        Which prices are outliers? The \$40 bike looks far below the rest --- is it an
        outlier? Justify each answer with a fence.
        \par\ansline{Only the \$900 bike: it is above the upper fence, \$270.}
        \par\ansline{The \$40 bike is not: \$40 is above the lower fence, \$30.}
        \par\vspace{4pt}
        Where does the right whisker of the modified boxplot stop? \ans{at \$200}

\end{enumerate}
\end{notesbox}

\boxguard[30]
\begin{notesbox}{Boxplots, Outliers, and Resistance, continued}
\begin{enumerate}[leftmargin=1.6em, itemsep=5pt, topsep=2pt, start=3]

  \item \textbf{With and without.} The shop sells the \$900 bike. Find the mean and the median
        before and after. (Sum of all eleven prices: \$2{,}310. Sum of the other ten:
        \$1{,}410.)

        \begin{center}
        \renewcommand{\arraystretch}{1.3}
        \begin{tabularx}{\linewidth}{@{} l Y Y @{}}
        \toprule
        & \textbf{With the \$900 bike (11)} & \textbf{Without it (10)} \\
        \midrule
        Mean   & \ans{\$2{,}310 $\div$ 11 $=$ \$210} & \ans{\$1{,}410 $\div$ 10 $=$ \$141} \\
        Median & \ans{\$150 (the 6th price)} & \ans{(\$140 $+$ \$150) $\div$ 2 $=$ \$145} \\
        \bottomrule
        \end{tabularx}
        \end{center}

  \item A classmate says, \emph{``Selling the \$900 bike pulls the median down just as much as
        the mean.''} Use item~3 to say whether she is right. Then name which summaries are
        resistant, using these too: with the \$900 bike, $\text{IQR} = \$60$ and
        $s \approx \$233$; without it, $\text{IQR} = \$60$ and $s \approx \$46$.
        \par\ansline{No --- the mean fell \$69 (\$210 to \$141); the median only \$5.}
        \par\ansline{Resistant: the median and IQR. Not resistant: the mean and SD.}

  \item Last season the shop sold $40$ used bikes. Read their cumulative relative frequency
        graph.

        \begin{center}
        \begin{tikzpicture}[x=0.03cm, y=0.034cm]
          \foreach \x in {50,100,...,350} { \draw[linegray!55] (\x,0) -- (\x,100); }
          \foreach \y in {5,10,...,100} { \draw[linegray!40] (50,\y) -- (350,\y); }
          \draw[thick] (50,0) -- (350,0);
          \draw[thick] (50,0) -- (50,100);
          \foreach \x in {50,100,...,350} { \node[below, font=\scriptsize] at (\x,0) {\$\x}; }
          \foreach \y in {0,25,50,75,100} { \node[left, font=\scriptsize] at (50,\y) {\y\%}; }
          \node[font=\scriptsize, anchor=north] at (200,-9) {price of the bike (dollars)};
          \node[font=\scriptsize, rotate=90, anchor=south] at (8,50) {cumulative percent};
          \draw[very thick, cerulean] (50,0) -- (100,15) -- (150,50) -- (200,75) -- (250,90) -- (300,95) -- (350,100);
          \foreach \x/\y in {50/0, 100/15, 150/50, 200/75, 250/90, 300/95, 350/100} {
            \fill[cerulean] (\x,\y) circle (2.2pt); }
        \end{tikzpicture}
        \end{center}

        Median price: \ans{\$150} \qquad Percent sold for \$250 or less: \ans{90\%}
        \par\vspace{4pt}
        \$200 is which percentile? Say what that means about the bikes.
        \par\ansline{The 75th --- 75\% of the 40 bikes sold for \$200 or less.}

  \item \emph{From Lesson 2.2:} the \textbf{range} is max $-$ min. Find it with and without the
        \$900 bike. Is the range resistant? Which spread should the shop's ad report, and why?

        \par\vspace{4pt}
        With: \ans{\$900 $-$ \$40 $=$ \$860} \qquad Without: \ans{\$200 $-$ \$40 $=$ \$160}
        \par\ansline{Not resistant --- it uses the \$900 itself. Report the IQR, \$60.}

\end{enumerate}
\end{notesbox}

\boxguard[5]   % the spiralbox needs about five baselines; a larger guard pushes it to a third page
\begin{spiralbox}
\textbf{Next lesson: Comparing Distributions.} Two boxplots on one axis, compared by center,
spread, and outliers --- with the resistant summaries you chose today.
\end{spiralbox}

\end{document}
